\section*{الفصل \ref{ch:b3:dna-repair} --- ثبات الجينوم: أذية DNA وإصلاحه وإعادة التركيب}

\begin{solution}{exo:b3:dna-repair:1}
اليوراسيل: عفوي (نزع أمين السيتوزين)، ويُصلَح \omterm{def:b3:dna-repair:excision}{بإصلاح استئصال القاعدة}
\omterm{def:b3:dna-repair:excision}{بغليكوزيلاز} اليوراسيل. وثنائي الثيمين: بيئي (فوق البنفسجي)، \omterm{def:b3:dna-repair:excision}{إصلاح
استئصال النوكليوتيد}. \omterm{def:b3:dna-repair:lesions}{و8-أوكسوغوانين}: عفوي (أكسدة)، \omterm{def:b3:dna-repair:excision}{إصلاح استئصال
القاعدة} \omterm{def:b3:dna-repair:excision}{بغليكوزيلاز} \omterm{def:b3:dna-repair:lesions}{8-أوكسوغوانين}. وعدم التوافق G--T خلف الشوكة:
خطأ تضاعف، \omterm{prop:b3:dna-repair:fidelity}{إصلاح عدم التوافق}. \omterm{def:b3:dna-repair:lesions}{والكسر المزدوج الطاق} في G1:
بيئي (إشعاع) أو عفوي؛ ويُصلَح \omterm{def:b3:dna-repair:dsb}{بوصل النهايات غير المتماثل}،
إذ لا صبيغية شقيقة متاحة.
\end{solution}

\begin{solution}{exo:b3:dna-repair:2}
يتعرف استئصال النوكليوتيد خاصية فيزيائية تشترك فيها الآفات الضخمة
كلها --- تشويه اللولب، أو بوليميراز RNA متوقف
--- ويزيل رقعة تحوي ما سبّبه. أما الآفات الصغيرة
(اليوراسيل \omterm{def:b3:dna-repair:lesions}{و8-أوكسوغوانين} والقواعد المؤلكلة) فلا تكاد تشوه اللولب،
فلا بد من تعرفها بكيميائها، \omterm{def:b3:dna-repair:excision}{بغليكوزيلاز} لكل نوع من القواعد
الخاطئة.
\end{solution}

\begin{solution}{exo:b3:dna-repair:3}
يتألف عدم التوافق من قاعدتين سويتين، والخاطئة منهما هي التي على
الطاق الجديد؛ فعلى المنظومة أن تعرف أي الطاقين جديد. وتعلّم
\emph{E.~coli} الطاقَ الوالدي بزمر ميثيل عند GATC،
وهي التي يفتقر إليها الطاق الجديد بضع دقائق؛ فيثلم MutH الطاق
غير الممثيل. وفي طافر \emph{dam} لا يُمثيل أي من الطاقين:
فيثلم MutH أيهما شاء، ويزيل الإصلاح القاعدة الصحيحة بقدر ما يزيل
الخاطئة --- فيثبّت نصف الأخطاء طفرات
(أي مطفِّر) --- وتصنع الثلمات على الطاقين عند مواقع GATC متجاورة
كسورًا مزدوجة الطاق.
\end{solution}

\begin{solution}{exo:b3:dna-repair:4}
نحو $35\times 2 = 70$ كسرًا مزدوج الطاق. وبعد ساعة يكون معظمها قد
وُصل: فبعمر نصف قرب \qty{30}{min}، تبقى نحو $70/4 \approx 18$
بؤرة. وبعد ست ساعات يكون كله تقريبًا قد أُصلح ($70/2^{12} < 1$)،
فلا يبقى إلا كسور قليلة معنّدة معقّدة.
\end{solution}

\begin{solution}{exo:b3:dna-repair:5}
(أ) $10^{-5}\times 10^{-2}\times 10^{-3} = 10^{-10}$؛ و$6.4\times
10^{9}\times 10^{-10} = 0.64$ طفرة في الانقسام. (ب) وبلا
\omterm{prop:b3:dna-repair:fidelity}{إصلاح عدم التوافق} $10^{-7}$؛ أي $640$ في الانقسام. (ج) وبلا تدقيق
$10^{-8}$؛ أي $64$ في الانقسام. فالطافر في \omterm{prop:b3:dna-repair:fidelity}{إصلاح عدم التوافق}
أخطر، بألف ضعف فوق السوي.
\end{solution}

\begin{solution}{exo:b3:dna-repair:6}
$k = \ln 2/\qty{0.25}{h} = \qty{2.8}{h^{-1}}$؛ و$D = 10^{4}/24 =
\qty{420}{h^{-1}}$؛ و$L^{*} = 420/2.8 \approx 150$ موضعًا بلا قاعدة
حاضرًا في أي لحظة. وبإبطاء الإصلاح عشرة أضعاف: $L^{*} \approx \num{1500}$.
وتمر الشوكة بكل موضع مرة، فتلقى الآفات الحاضرة هناك في
تلك اللحظة؛ وبالمتوسط على الجينوم تلقى نحو $L^{*}$ منها
--- نحو $150$ في السوي، و$\num{1500}$ مع الدواء.
\end{solution}

\begin{solution}{exo:b3:dna-repair:7}
$k_{\text{NHEJ}} = \ln 2/0.5 = \qty{1.39}{h^{-1}}$، و$k_{\text{HR}} =
\ln 2/4 = \qty{0.17}{h^{-1}}$؛ ونسبة إعادة التركيب $0.17/1.56 = 0.11$؛
وعمر النصف $\ln 2/1.56 = \qty{27}{min}$. والشوكة المنهارة تعطي
كسرًا \emph{بنهاية واحدة}: فلا نهاية ثانية يصلها وصل النهايات،
ولا يعيد الشوكة إلا إعادة التركيب مع الشقيقة --- وهي على كل حال
المسلك الدقيق الوحيد.
\end{solution}

\begin{solution}{exo:b3:dna-repair:8}
(أ) لينش: سواتل مكروية غير مستقرة (تختلف أطوالها بين خلايا
الورم)، وسرطانات القولون والمستقيم وبطانة الرحم والمعدة والمبيض،
واستجابة سوية لضوء الشمس. (ب) XP-A: سواتل مكروية مستقرة، وحروق شمس
شديدة عند أدنى تعرض، وسرطانات جلد وعين في الطفولة، وفي
هذه الزمرة تنكّس عصبي تدريجي. (ج) XP-V:
سواتل مكروية مستقرة، وإصلاح الاستئصال سوي فالحرق أخف،
لكن الثنائيات التي تفلت من الاستئصال تُتجاوز ببوليميرازات كثيرة
الخطأ --- فسرطانات جلد، أمتنُ عمرًا من زمرة~A.
\end{solution}

\begin{solution}{exo:b3:dna-repair:9}
خلايا الكبد: استأصل Cre الإكسون~3 من الأليلين معًا (حلقةً
تُفقد)، فتكون المورّثة معطّلة في كل خلية كبدية منذ الولادة
--- أي حذف خاص بالكبد. والعصبونات: لا Cre فيها، فالأليلان
المؤطَّران سليمان عاملان. وفي الاتجاهين المتعاكسين: يقلب Cre
الإكسون ويعيد قلبه بلا نهاية، فيحمل نحو نصف الخلايا في أي وقت
الإكسونَ المقلوب غير العامل --- أي فسيفساء لا حذفًا نظيفًا.
\end{solution}

\begin{solution}{exo:b3:dna-repair:10}
كلما شُطر LexA هبط تركيزه تدريجًا. والمشغّلات التي ترتبط به
ارتباطًا ضعيفًا تُخلى عند هبوط صغير --- \emph{uvrA}،
و\emph{uvrB}، و\emph{recA} --- فيبدأ الإصلاح الدقيق أولًا؛ أما
مشغّلات \emph{sulA} والبوليميرازات كثيرة الخطأ فترتبط به
ارتباطًا شديدًا ولا تُخلى إلا حين يكاد يفنى، أي حين
تستمر الأذية زمنًا طويلًا. والبكتيريا بلا البوليميرازات
كثيرة الخطأ تموت عند شوكات تحجبها آفات لا تستطيع
تجاوزها، ومن غير تطفير محرَّض تتطور مقاومتها للمضاد
الحيوي أبطأ بكثير --- وهو عيب للبكتيريا ومكسب
للمريض؛ ولهذا تُدرس مثبطات \omterm{prop:b3:dna-repair:sos}{استجابة SOS}.
\end{solution}

\begin{solution}{exo:b3:dna-repair:11}
السلسلة: تثبيط PARP $\to$ بقاء الكسور المفردة الطاق $\to$ تحويل
شوكة تضاعف كلًا منها إلى \omterm{def:b3:dna-repair:lesions}{كسر مزدوج الطاق} بنهاية واحدة
$\to$ اقتضاء الإصلاح إعادةَ تركيب متماثلة $\to$ عجز خلية الورم،
وهي معوزة \omterm{def:b3:dna-repair:dsb}{BRCA}، عن ذلك $\to$ وصل خاطئ، وفقد صبغيات، وموت.
أما الخلايا السوية، ولها أليل عامل واحد، فتعيد التركيب وتنجو.
والمقاومة: (1) طفرة ثانية في \emph{\omterm{def:b3:dna-repair:dsb}{BRCA}} تعيد إطار القراءة
وبروتينًا عاملًا --- وتُكشف بتسلسل الورم أو
DNA الورمي الجائل في الدم، وبعودة بؤر \omterm{def:b3:dna-repair:dsb}{Rad51} بعد
التشعيع؛ (2) فقد البروتينات التي تحجب قطع النهايات
(53BP1--shieldin)، وهو يعيد جزءًا من إعادة التركيب بلا BRCA1
--- ويُكشف بغيابها في التلوين وبعودة بؤر \omterm{def:b3:dna-repair:dsb}{Rad51}؛
وكذلك طفرات PARP1 التي تمنع احتباسه على DNA، أو مضخات
النَّفْث الدوائي، وتُكشف بالتسلسل وبالتعبير تباعًا.
\end{solution}

\begin{solution}{exo:b3:dna-repair:12}
لولبان متماثلان، أحدهما يحمل $A$ والآخر $a$، يتبادلان
طاقات مفردة عبر الواسم: فيصير لكل منهما طاق من كل
والد على امتداد المسار --- أي ازدواجية متغايرة بعدم توافق $A/a$. فإذا
حوّل \omterm{prop:b3:dna-repair:fidelity}{إصلاح عدم التوافق} عدمَ التوافق على اللولب المغزوّ إلى $A$،
حملت ثلاث صبيغيات $A$ وواحدة $a$: أي $3{:}1$؛ وإذا حُوّل إلى $a$:
كان $1{:}3$؛ وإذا تُرك بلا إصلاح، انعزل $A$ وَ$a$ في الصبيغية عند
أول انقسام فتيلي بعد الانقسام المنصّف، فأعطى بوغين بنتين مختلفين
(انعزال بعد الانقسام المنصّف، أي $5{:}3$ في كيس ثماني الأبواغ).
ومسار الازدواجية المتغايرة بضع مئات إلى بضعة آلاف من أزواج القواعد
حول الكسر البادئ، فلا يُحوَّل إلا الواسمات الواقعة فيه، وهي
بحكم البناء مجاورة للتصالب.
\end{solution}

\begin{solution}{pb:b3:dna-repair:1}
\textbf{1.} $10^{4} + 400 + 10^{4} \approx 2\times 10^{4}$ آفة في
اليوم، و$7\times 10^{6}$ في السنة --- أي نحو واحدة لكل كيلوقاعدة من
الجينوم في السنة.
\textbf{2.} $10^{-5}\times 10^{-2}\times 10^{-3} = 10^{-10}$ لكل زوج قواعد؛
و$6.4\times 10^{9}\times 10^{-10} = 0.64$ طفرة في الانقسام.
\textbf{3.} $0.64\times 0.015 \approx 0.01$ طفرة مشفِّرة في
الانقسام؛ وعلى مدى $50$ انقسامًا، نحو $0.5$.
\textbf{4.} $10^{-7}$ لكل زوج قواعد، و$640$ طفرة في الانقسام. والخلايا
المنزلقة: بلا إصلاح $10^{-4}\times 40\times 10^{9} = 4\times 10^{6}$؛
ومعه $10^{-7}\times 40\times 10^{9} = 4000$ --- أي فرق بألف ضعف،
يظهر بوصفه \omterm{def:b3:dna-repair:mmr}{عدم استقرار السواتل المكروية}.
\textbf{5.} اليوراسيل غريب عن DNA، ويتعرفه \omterm{def:b3:dna-repair:excision}{غليكوزيلاز} اليوراسيل
ويستأصله: فيُصلَح. أما الثيمين الآتي من 5-ميثيل سيتوزين فقاعدة
سوية قبالة G؛ ولا شيء يعلّمه خطأً خارج التضاعف،
فيصير طفرة C$\to$T --- وهي الآلية التي استنزفت
\omterm{def:b3:chromatin-epigenetics:methylation}{CpG} من الجينوم.
\textbf{6.} يزيل \omterm{def:b3:dna-repair:excision}{غليكوزيلاز} \omterm{def:b3:dna-repair:lesions}{8-أوكسوغوانين} القاعدةَ المؤكسدة ما دامت
قبالة C (قبل التضاعف)؛ ويزيل \omterm{def:b3:dna-repair:excision}{الغليكوزيلاز} الثاني
الأدنينَ المدخل خطأً قبالة \omterm{def:b3:dna-repair:lesions}{8-أوكسوغوانين} (بعد التضاعف، فيعيد
فرصة إصلاح G)؛ وتهدم الحلمهة dGTP المؤكسد فلا
يُدخَل قبالة A. ومتى أخفقت كلها، ازدوج \omterm{def:b3:dna-repair:lesions}{8-أوكسوغوانين} مع A وثبّتت
الجولة التالية تبدّلًا عرضيًا G$\cdot$C$\to$T$\cdot$A.
\textbf{7.} في السوي $k = \ln 2/2 = \qty{0.35}{h^{-1}}$؛ وفي جفاف الجلد $k = \ln 2/70
= \qty{0.0099}{h^{-1}}$.
\textbf{8.} في السوي $10^{5} e^{-2.77} \approx 6200$؛ وفي جفاف الجلد $10^{5}
e^{-0.079} \approx \num{92000}$.
\textbf{9.} الطفرات $6.2$ و$92$ في الخلية؛ والمشفِّرة $0.09$ و$1.4$.
\textbf{10.} عند $500$ تعرض: $47$ طفرة مشفِّرة (في السوي) و$690$
(في جفاف الجلد)، مقابل $0.5$ من التضاعف --- فضوء الشمس يهيمن على
حمل الطفرات في الجلد حتى عند الشخص السوي.
\textbf{11.} التسلسل المشفِّر $9.6\times 10^{7}$ زوج قواعد؛ والمورّثات
القائدة منه $4\times 10^{4}/9.6\times 10^{7} = 4.2\times 10^{-4}$. ومتوسط
الإصابات $690\times 4.2\times 10^{-4} = 0.29$؛ و$P(\ge 1) = 1 - e^{-0.29} =
0.25$. ومن $10^{9}$ خلية، تحمل $2.5\times 10^{8}$ طفرة قائدة
(وفي الطفل السوي: المتوسط $0.02$، أي \qty{2}{\%}).
\textbf{12.} الثنائي ليس طفرة؛ ولا يصير طفرة إلا حين
تنسخه شوكة بتجاوز كثير الخطأ، فالخلايا المنقسمة ذات الثنائيات
الكثيرة غير المصلَحة هي التي تطفر. والضوء فوق البنفسجي لا يبلغ
إلا الجلد والعين؛ أما الأنسجة الداخلية فلا تتلقى ثنائيات، وهي في
جفاف الجلد شبه سوية.
\textbf{13.} $70$ كسرًا؛ وبعد \qty{1}{h} بوصل النهايات وحده، $70\times
2^{-2} \approx 18$ بؤرة.
\textbf{14.} $k_{\text{NHEJ}} = \qty{1.39}{h^{-1}}$، و$k_{\text{HR}} =
\qty{0.17}{h^{-1}}$؛ ونسبة إعادة التركيب $0.17/1.56 = 0.11$.
\textbf{15.} عمر النصف $\ln 2/1.56 = \qty{27}{min}$؛ وبعد \qty{1}{h}
يبقى $70\,e^{-1.56} \approx 15$ كسرًا.
\textbf{16.} $70\times 0.015\times 0.7 \approx 0.7$ مورّثة معطَّلة في
الخلية.
\textbf{17.} عند $n = 70$: $2415$ زوجًا $\times 3\times 10^{-4} = 0.72$
انتقال. وعند $n = 7$: $21\times 3\times 10^{-4} = 0.006$. فعدد الكسور
$\propto D$ لكن عدد \emph{أزواج} الكسور المتزامنة
$\propto n^{2} \propto D^{2}$.
\textbf{18.} $D = \qty{7}{h^{-1}}$، و$L^{*} = 7/1.39 \approx 5$ كسور
حاضرة في اللحظة: أي $10$ أزواج، و$0.003$ انتقال --- أي أقل مئتي مرة
من الجرعة نفسها في ومضة. والتجزئة تتيح للنسيج السوي أن يصلح بين
الجرعات وتُبقي الكسور المتزامنة قليلة.
\textbf{19.} بإصلاح سليم: $e^{-2/1.5} = 0.26$ عند \qty{2}{Gy}،
و$e^{-4} = 0.018$ عند \qty{6}{Gy}. وبلا وصل نهايات: $e^{-4} = 0.018$ و
$e^{-12} = 6\times 10^{-6}$. و$D_{0}$ هي الجرعة التي تترك وسطيًا
آفة قاتلة واحدة غير مصلَحة في الخلية (فتكون النجاة $1/e$).
\textbf{20.} إن \qty{11}{\%} من كسور G2 ذات النهايتين التي كانت إعادة
التركيب لتصلحها إصلاحًا دقيقًا تمر الآن بوصل النهايات، وكل كسر شوكة
بنهاية واحدة --- وهو ما لا يصلحه وصل النهايات إصلاحًا صحيحًا
البتة --- يوصَل خطأً: فيتراكم في الجينوم على مر السنين حذوف
وانتقالات وتغيرات في عدد النسخ، وهي ندوب أورام \emph{\omterm{def:b3:dna-repair:dsb}{BRCA}}
المميزة.
\textbf{21.} $10^{4}\times 0.01 = 100$ كسر بنهاية واحدة في اليوم.
والكسر بنهاية واحدة لا شريك له؛ فلا يملك وصل النهايات إلا أن
يتركه أو يصله بنهاية لا صلة لها به، فينتج انتقال أو فقد.
\textbf{22.} أليل \emph{BRCA2} عامل واحد يصنع بروتينًا يكفي
لإعادة التركيب: فتصلح الخلايا السوية كسور شوكاتها وتنجو من
المثبط.
\textbf{23.} إزاحة إطار ثانية بعد الأولى تعيد إطار
القراءة فتعطي بروتين BRCA2 مقصَّرًا لكنه عامل جزئيًا، فتعود
إعادة التركيب ولم يعد الدواء يقتل. ويُكشف ذلك بتسلسل
\emph{BRCA2} في الورم المقاوم أو في DNA الورمي الجائل،
ووظيفيًا بعودة ظهور بؤر \omterm{def:b3:dna-repair:dsb}{Rad51} بعد التشعيع.
\textbf{24.} يحتاج القتل إلى منظومة \omterm{prop:b3:dna-repair:fidelity}{إصلاح عدم التوافق} لتشتبك مع
القواعد المؤلكلة؛ وورم لينش الذي تنقصه متسامح --- فيخفق الدواء في
قتله (وتُعالج تلك الأورام بوسائل أخرى، منها المعالجة المناعية،
التي يجعلها حملُها الطفري الثقيل حساسةً لها).
\textbf{25.} نحو \num{92000} ثنائي عند الشوكة في خلية جفاف الجلد؛
و$10^{-10}$ طفرة لكل زوج قواعد في الانقسام؛ وإعادة التركيب تصلح
\qty{11}{\%} من كسور G2.
\end{solution}
